\documentclass[11pt,a4paper]{article}
\usepackage[utf8]{inputenc}
\usepackage[vietnamese]{babel}
\usepackage{amsmath,amssymb}
\usepackage{geometry}
\geometry{a4paper, total={180mm,260mm}, left=15mm, right=15mm, top=18mm, bottom=18mm}
\usepackage{tikz}
\usetikzlibrary{shapes.geometric, arrows.meta, positioning, calc}
\usepackage{tcolorbox}
\tcbuselibrary{skins,breakable}
\usepackage{xcolor}
\usepackage{multicol}
\usepackage{array}
\usepackage{fontawesome}
\usepackage{fancyhdr}

\pagestyle{fancy}
\fancyhf{}
\renewcommand{\headrulewidth}{0.5pt}
\renewcommand{\footrulewidth}{0.5pt}
\fancyhead[L]{\small\bfseries\color{teal}\faBook\ SỔ TAY KIẾN THỨC TOÁN 3}
\fancyhead[R]{\small\bfseries\color{teal}Lớp Toán Thầy Lĩnh -- loptoanthaylinh.vn}
\fancyfoot[C]{\small\color{gray}-- Sổ tay kiến thức: Bài 55. Phép trừ trong phạm vi 10 000 --}

\definecolor{primary}{HTML}{005F73}
\definecolor{secondary}{HTML}{0A9396}
\definecolor{accent}{HTML}{EE9B00}
\definecolor{danger}{HTML}{AE2012}
\definecolor{boxbg}{HTML}{F4F6F8}

\begin{document}

% --- BANNER TIÊU ĐỀ ---
\begin{tcolorbox}[
  enhanced,
  colback=white,
  colframe=primary,
  arc=3mm,
  boxrule=1.2pt,
  left=4mm, right=4mm, top=2.5mm, bottom=2.5mm,
  drop shadow=black!10!white
]
  \centering
  {\Large\bfseries\color{primary}\faBookmark\ BÀI 55: PHÉP TRỪ TRONG PHẠM VI 10 000}\\[1.5mm]
  {\small\itshape\color{black!75}Kĩ thuật đặt tính, trừ từ phải sang trái, mượn 1 chục và trả nhớ chính xác trên 1 trang A4}
\end{tcolorbox}

\vspace{-2mm}


\noindent
\begin{minipage}[t]{0.49\textwidth}\vspace{0pt}
\begin{tcolorbox}[
  enhanced,
  colback=blue!4!white,
  colframe=blue!75!black,
  arc=2mm, boxrule=0.8pt,
  left=2mm, right=2mm, top=1.8mm, bottom=1.8mm,
  title={\footnotesize\bfseries\color{white}\faKey\ 1. Quy trình Đặt tính và Trừ}
]
{\footnotesize
\begin{itemize}\setlength{\itemsep}{1.5pt}\setlength{\leftskip}{-2.5mm}
  \item \textbf{Đặt tính:} Viết số trừ dưới số bị chia sao cho các chữ số ở cùng một hàng thẳng cột. Viết dấu $-$, kẻ vạch ngang.
  \item \textbf{Thứ tự trừ:} Thực hiện trừ từ \textbf{phải sang trái}:
  $$\text{Đơn vị} \to \text{Chục} \to \text{Trăm} \to \text{Nghìn}$$
  \item \textbf{Quy tắc có mượn (có nhớ):}
  Nếu chữ số số bị trừ nhỏ hơn số trừ, mượn $1$ chục ở hàng liền trước (cộng thêm $10$) để trừ, rồi \textbf{nhớ $1$ trả sang hàng tiếp theo} của số trừ.
\end{itemize}
}
\end{tcolorbox}
\end{minipage}
\hfill
\begin{minipage}[t]{0.49\textwidth}\vspace{0pt}
\begin{tcolorbox}[
  enhanced,
  colback=teal!4!white,
  colframe=teal!75!black,
  arc=2mm, boxrule=0.8pt,
  left=2mm, right=2mm, top=1.8mm, bottom=1.8mm,
  title={\footnotesize\bfseries\color{white}\faCheckSquare\ 2. Kĩ năng thử lại phép trừ}
]
{\footnotesize
\begin{itemize}\setlength{\itemsep}{1.5pt}\setlength{\leftskip}{-2.5mm}
  \item \textbf{Thử lại bằng phép cộng:}
  $$\text{Hiệu} + \text{Số trừ} = \text{Số bị trừ}$$
  Nếu kết quả cộng lại đúng bằng số bị trừ thì phép trừ đã làm chính xác $100\%$.
  \item \textbf{Tìm thành phần chưa biết:}
  $$\text{Số bị trừ} = \text{Hiệu} + \text{Số trừ}$$
  $$\text{Số trừ} = \text{Số bị trừ} - \text{Hiệu}$$
\end{itemize}
}
\end{tcolorbox}
\end{minipage}

\vspace{-2.5mm}
% --- KHỐI 3: SƠ ĐỒ TƯ DUY TIKZ ---
\noindent
\begin{tcolorbox}[
  enhanced,
  colback=orange!4!white,
  colframe=orange!80!black,
  arc=2mm, boxrule=0.8pt,
  left=2mm, right=2mm, top=1.5mm, bottom=1.5mm,
  title={\footnotesize\bfseries\color{white}\faSitemap\ SƠ ĐỒ TƯ DUY: NGUYÊN LÝ MƯỢN VÀ TRẢ KHI TRỪ CÓ NHỚ}
]
\centering
\begin{tikzpicture}[
  box/.style={draw, rounded corners=2pt, align=center, font=\scriptsize\bfseries, inner sep=2.5pt, minimum height=0.75cm},
  arr/.style={-{Stealth[length=2mm]}, thick}
]
  \node[box, fill=purple!15, draw=purple!80!black, text width=2.8cm] (t1) at (0, 0) {CHỮ SỐ BÉ HƠN\\Không trừ được\\$a < b$};
  \node[box, fill=teal!15, draw=teal!80!black, text width=2.8cm] (t2) at (4.0, 0) {MƯỢN $1$ CHỤC\\$(10 + a) - b$\\Viết kết quả};
  \node[box, fill=blue!12, draw=blue!80!black, text width=2.8cm] (t3) at (8.0, 0) {TRẢ NHỚ $1$\\Thêm $1$ vào chữ số\\hàng sau của số trừ};
  \node[box, fill=orange!15, draw=orange!80!black, text width=2.2cm] (t4) at (11.4, 0) {TRỪ TIẾP\\Hàng liền kề};

  \draw[arr, purple!80!black] (t1) -- (t2);
  \draw[arr, teal!80!black] (t2) -- (t3);
  \draw[arr, blue!80!black] (t3) -- (t4);
\end{tikzpicture}
\end{tcolorbox}

\vspace{-2.5mm}
% --- KHỐI 4: VÍ DỤ & MẸO TRÁNH BẪY ---
\noindent
\begin{minipage}[t]{0.49\textwidth}\vspace{0pt}
\begin{tcolorbox}[
  enhanced,
  colback=green!4!white,
  colframe=green!60!black,
  arc=2mm, boxrule=0.8pt,
  left=2mm, right=2mm, top=1.8mm, bottom=1.8mm,
  title={\footnotesize\bfseries\color{white}\faCheckCircle\ Ví dụ thực hành chuẩn mực}
]
{\footnotesize
\textbf{1. Đặt tính trừ có nhớ:}\\
$\begin{array}{r@{}l}
  7245 & \\
  -\quad 3618 & \\
  \hline
  \mathbf{3627} &
\end{array}$
\qquad
\begin{minipage}[c]{0.5\linewidth}
-- $5$ không trừ được $8$, lấy $15 - 8 = 7$, viết $7$ nhớ $1$.\\
-- $1$ thêm $1$ là $2$; $4 - 2 = 2$, viết $2$.\\
-- $2$ không trừ được $6$, lấy $12 - 6 = 6$, viết $6$ nhớ $1$.\\
-- $3$ thêm $1$ là $4$; $7 - 4 = 3$, viết $3$.
\end{minipage}

\textbf{2. Trừ qua chữ số 0:}\\
$5000 - 2430 = \mathbf{2570}$.
}
\end{tcolorbox}
\end{minipage}
\hfill
\begin{minipage}[t]{0.49\textwidth}\vspace{0pt}
\begin{tcolorbox}[
  enhanced,
  colback=red!4!white,
  colframe=red!70!black,
  arc=2mm, boxrule=0.8pt,
  left=2mm, right=2mm, top=1.8mm, bottom=1.8mm,
  title={\footnotesize\bfseries\color{white}\faExclamationTriangle\ Mẹo tránh bẫy trắc nghiệm}
]
{\footnotesize
\begin{itemize}\setlength{\itemsep}{1.5pt}\setlength{\leftskip}{-2.5mm}
  \item \textbf{Bẫy quên trả số nhớ khi trừ:}\\
  Mượn $1$ ở hàng đơn vị nhưng sang hàng chục quên thêm $1$ vào số trừ $\implies$ kết quả bị lớn hơn $1$ chục.
  \item \textbf{Bẫy trừ ngược từ dưới lên trên:}\\
  Thấy $5 < 8$, nhiều bạn lấy $8 - 5 = 3$ là \textbf{SAI}! Phải luôn lấy số trên trừ số dưới, nếu không đủ thì phải mượn.
\end{itemize}
}
\end{tcolorbox}
\end{minipage}


\end{document}
